\documentclass[10pt, letterpaper]{article}

% Packages:
\usepackage[
    ignoreheadfoot, % set margins without considering header and footer
    top=1.5 cm, % seperation between body and page edge from the top
    bottom=1.5 cm, % seperation between body and page edge from the bottom
    left=1.5 cm, % seperation between body and page edge from the left
    right=1.5 cm, % seperation between body and page edge from the right
    footskip=1.0 cm, % seperation between body and footer
    % showframe % for debugging 
]{geometry} % for adjusting page geometry
\usepackage{titlesec} % for customizing section titles
\usepackage{tabularx} % for making tables with fixed width columns
\usepackage{array} % tabularx requires this
\usepackage[dvipsnames]{xcolor} % for coloring text
\definecolor{primaryColor}{RGB}{0, 0, 0} % define primary color
\usepackage{enumitem} % for customizing lists
\usepackage{fontawesome5} % for using icons
\usepackage{amsmath} % for math
\usepackage{ragged2e}   % provides \justifying
\usepackage{microtype}  % improves justification/spacing (recommended)
\usepackage[
    pdftitle={John Doe's CV},
    pdfauthor={John Doe},
    pdfcreator={LaTeX with RenderCV},
    colorlinks=true,
    urlcolor=primaryColor
]{hyperref} % for links, metadata and bookmarks
\usepackage[pscoord]{eso-pic} % for floating text on the page
\usepackage{calc} % for calculating lengths
\usepackage{bookmark} % for bookmarks
\usepackage{lastpage} % for getting the total number of pages
\usepackage{changepage} % for one column entries (adjustwidth environment)
\usepackage{paracol} % for two and three column entries
\usepackage{ifthen} % for conditional statements
\usepackage{needspace} % for avoiding page brake right after the section title
\usepackage{iftex} % check if engine is pdflatex, xetex or luatex

% Ensure that generate pdf is machine readable/ATS parsable:
\ifPDFTeX
    \input{glyphtounicode}
    \pdfgentounicode=1
    \usepackage[T1]{fontenc}
    \usepackage[utf8]{inputenc}
    \usepackage{lmodern}
\fi

\usepackage{charter}

% Some settings:
% \raggedright
\AtBeginEnvironment{adjustwidth}{\partopsep0pt} % remove space before adjustwidth environment
\pagestyle{empty} % no header or footer
\setcounter{secnumdepth}{0} % no section numbering
\setlength{\parindent}{0pt} % no indentation
\setlength{\topskip}{0pt} % no top skip
\setlength{\columnsep}{0.15cm} % set column seperation
\pagenumbering{gobble} % no page numbering

\titleformat{\section}{\needspace{4\baselineskip}\bfseries\large}{}{0pt}{}[\vspace{1pt}\titlerule]

\titlespacing{\section}{
    % left space:
    -1pt
}{
    % top space:
    0.3 cm
}{
    % bottom space:
    0.2 cm
} % section title spacing

\renewcommand\labelitemi{$\vcenter{\hbox{\small$\bullet$}}$} % custom bullet points
\newenvironment{highlights}{
    \begin{itemize}[
        topsep=0.10 cm,
        parsep=0.10 cm,
        partopsep=0pt,
        itemsep=0pt,
        leftmargin=0 cm + 10pt
    ]
}{
    \end{itemize}
} % new environment for highlights


\newenvironment{highlightsforbulletentries}{
    \begin{itemize}[
        topsep=0.10 cm,
        parsep=0.10 cm,
        partopsep=0pt,
        itemsep=0pt,
        leftmargin=10pt
    ]
}{
    \end{itemize}
} % new environment for highlights for bullet entries

\newenvironment{onecolentry}{
    \begin{adjustwidth}{
        0 cm + 0.00001 cm
    }{
        0 cm + 0.00001 cm
    }
}{
    \end{adjustwidth}
} % new environment for one column entries

\newenvironment{twocolentry}[2][]{
    \onecolentry
    \def\secondColumn{#2}
    \setcolumnwidth{\fill, 4.5 cm}
    \begin{paracol}{2}
}{
    \switchcolumn \raggedleft \secondColumn
    \end{paracol}
    \endonecolentry
} % new environment for two column entries

\newenvironment{threecolentry}[3][]{
    \onecolentry
    \def\thirdColumn{#3}
    \setcolumnwidth{, \fill, 4.5 cm}
    \begin{paracol}{3}
    {\raggedright #2} \switchcolumn
}{
    \switchcolumn \raggedleft \thirdColumn
    \end{paracol}
    \endonecolentry
} % new environment for three column entries

\newenvironment{header}{
    \setlength{\topsep}{0pt}\par\kern\topsep\centering\linespread{1.5}
}{
    \par\kern\topsep
} % new environment for the header

\newcommand{\placelastupdatedtext}{% \placetextbox{<horizontal pos>}{<vertical pos>}{<stuff>}
  \AddToShipoutPictureFG*{% Add <stuff> to current page foreground
    \put(
        \LenToUnit{\paperwidth-2 cm-0 cm+0.05cm},
        \LenToUnit{\paperheight-1.0 cm}
    ){\vtop{{\null}\makebox[0pt][c]{
        \small\color{gray}\textit{Last updated in September 2024}\hspace{\widthof{Last updated in September 2024}}
    }}}%
  }%
}%

\newcommand{\sectionwithkeywords}[3]{%
  \section{\texorpdfstring{%
    \begin{tabular*}{\linewidth}{@{}l@{\extracolsep{\fill}}r@{}}%
      #1 & {\normalsize #2}%
    \end{tabular*}%
  }{#3}}%
}



% save the original href command in a new command:
\let\hrefWithoutArrow\href

% new command for external links:



\begin{document}\small
    \newcommand{\AND}{\unskip
        \cleaders\copy\ANDbox\hskip\wd\ANDbox
        \ignorespaces
    }
    \newsavebox\ANDbox
    \sbox\ANDbox{$|$}

    \begin{header}
        {\fontsize{22pt}{22pt}\bfseries Xingyi Zhao}\hfill
        \hrefWithoutArrow{https://github.com/xingyizhao}{xingyi.zhao.github} \,\textbullet\, 
        \hrefWithoutArrow{tel:+14352945827}{435-294-5827} \,\textbullet\,
        \hrefWithoutArrow{mailto:xingyi.zhao@usu.edu}{xingyi.zhao@usu.edu}
    \end{header}

    \vspace{8 pt}

    \sectionwithkeywords
      {Research Interests}
      {LLM Security \textbullet\ Faithful Reasoning \textbullet\ Safe Post-Training (SFT, RLHF) \textbullet\ LLM Agent}
      {Research Interests}


    \section{Education}
    \begin{highlights}
    
        \item\begin{twocolentry}{Sep. 2021 - Present}
            \textbf{Utah State University, UT, USA}\\
            \textbf{Ph.D} in Computer Science (GPA: 4.0/4.0)
            \end{twocolentry}

        \item\begin{twocolentry}{Sep. 2016 – June 2019}
            \textbf{Huazhong University of Science and Technology, China}\\
            \textbf{M.Eng.} in Control Science (GPA: 86.02/100)
            \end{twocolentry}
        
        \item\begin{twocolentry}{Sep. 2012 – June 2016}
            \textbf{Huazhong University of Science and Technology, China}\\
            \textbf{B.Eng.} in Measurement and Control Technology and Instrument (GPA: 3.58/4.0)
            \end{twocolentry}
    \end{highlights}


    \section{Research Experience}
        \begin{twocolentry}{
            Sept 2021 - Present 
        }
            \textbf{Research Assistant}, Department of Computer Science, Utah State University\end{twocolentry}
        \vspace{0.10 cm}
        \begin{onecolentry}
            \begin{highlights}
                \item \textbf{Investigating LLM Backdoor Collapse and Making Backdoors Persistent.} Observed that LLM backdoors often collapse after trigger-free downstream SFT because the backdoor optimum lies in a sharp, narrow basin, so small parameter drift can cause backdoor forgetting. Proposed a novel algorithm, BAD-BOOM, which smooths backdoor-sensitive parameters and moves the backdoor optimum into a broader, smoother basin. BAD-BOOM can maintain ASR $\geq$90\%, whereas AdamW often drops to 0\% after downstream SFT.

                \item \textbf{Developing LLM Agents for Virtual Bio-Lab.} Built an agent workflow for neural cell drug-response analysis. Designed harness engineering modules that package electrophysiology signals and derived summaries (e.g., firing-rate) into structured prompts, enabling the agent to classify drug-stimulated vs. non-stimulated cells. Implemented reasoning capture (CoT-style rationales) and evidence references to improve interpretability, supporting downstream analysis for hypothesis generation in drug screening.

                \item \textbf{Designing Robust Gradient Ascent for LLM Malicious Behavior Unlearning.} Identified a key limitation of vanilla gradient-ascent unlearning: GA keeps maximizing poisoned-sample loss with no natural stopping point, leading to over-unlearning and trigger shifting. Proposed Robust Gradient Ascent (RGA): adaptively decays GA strength using a KL-divergence penalty between the current policy and a reference policy on poisoned samples, preventing loss from growing indefinitely. RGA removes malicious behavior with performance comparable to clean retraining, while vanilla GA causes severe trigger shifting ($\approx$50\% accuracy on poisoned tests). 

                \item \textbf{Securing Backdoor Defensive Algorithm.} Recognized a real-world threat that end users may download a poisoned PLM (e.g., from HuggingFace) without access to backdoor knowledge like trigger type. Proposed PURE—(1) prune low-utility/trigger-hijacked attention heads in BERT-like models (BERT, DistilBERT) identified via low [CLS] token attention-variance on clean data, and (2) add an attention-normalization regularizer to prevent [CLS] token from over-focusing on specific tokens—reducing attack success from $\geq$90\% to as low as 7.99, while largely preserving clean accuracy.

                \item \textbf{Advancing Efficient Adversarial Attacks on Text Classifiers}. Identified word-level attacks face a trade-off—combinatorial search (e.g., GA/PSO) is slow, while greedy methods are faster but often yield lower-quality/less effective adversarial examples. Proposed TAMPERS, a two-stage word-level attack that combines greedy vulnerable-word selection with genetic search in a reduced space to produce semantically preserved adversarial edits. Achieves higher success with fewer changes and runs faster than GA-based attacks.        
            \end{highlights}
        \end{onecolentry}


    \section{Selected Publications (\href{https://scholar.google.com/citations?user=yVq7_1sAAAAJ&hl=en}{Google Scholar})}

    \begin{highlights}

        \item\begin{onecolentry}
            \mbox{\textbf{\underline{Xingyi Zhao}}}, \mbox{Shuhan Yuan}, et al. Broadening the Backdoor Basin: Understanding LLM Backdoors Collapse and Making Backdoors Persistent. In Proceedings of the 43rd International Conference on Machine Learning: ICML 2026.
        \end{onecolentry}
        
        \item\begin{onecolentry}
            \mbox{\textbf{\underline{Xingyi Zhao}}}, \mbox{Shuhan Yuan}, et al. Don't Shift the Trigger: Robust Gradient Ascent for Backdoor Unlearning. In Proceedings of the 14th International Conference on Learning Representation: ICLR 2026.
        \end{onecolentry}

        \item\begin{onecolentry}
            \mbox{\textbf{\underline{Xingyi Zhao}}}, \mbox{Shuhan Yuan}, et al. Defense against Backdoor Attack on Pre-trained Language Models via Head Pruning and Attention Normalization. In Proceedings of the 41st International Conference on Machine Learning: ICML 2024. 
            \end{onecolentry}
    
        \item\begin{onecolentry}
            \mbox{\textbf{\underline{Xingyi Zhao}}}, \mbox{Shuhan Yuan}, et al. Generating Textual Adversaries with Minimal Perturbation. In Findings of the Association for Computational Linguistics: EMNLP 2022.
            \end{onecolentry}

    \end{highlights}

    \section{Technical Skills}
    \begin{onecolentry}
    \begin{highlights}
      \item \textbf{Programming Languages:} Python, CUDA, C/C++, Java, SQL, MATLAB
      \item \textbf{ML / Deep Learning:} NumPy, PySpark, Pandas, scikit-learn, PyTorch, TensorFlow, Hugging Face Transformers
      \item \textbf{LLM Training:} TRL, VERL, LoRA/PEFT, DeepSpeed, Accelerate; (Models: LLaMA, Qwen, GPT-OSS)
      \item \textbf{NLP / Reasoning:} Post-Training (SFT, PPO, DPO, GRPO), CoT, Harness Engineering
      \item \textbf{Agent Tools:} LangChain, CrewAI, AutoGPT
    \end{highlights}
    \end{onecolentry}

       
\end{document}